% OpenStax Calculus Volume 1 -- Section 4.6 Exercises: Limits at Infinity and Asymptotes
% Exercises reproduced from OpenStax, licensed under CC BY 4.0.
% Source: https://openstax.org/books/calculus-volume-1/pages/4-6-limits-at-infinity-and-asymptotes
\documentclass{exercises}
\title{OpenStax Calculus Volume 1}
\subtitle{Section 4.6 Exercises: Limits at Infinity and Asymptotes}
\sourceurl{https://openstax.org/books/calculus-volume-1/pages/4-6-limits-at-infinity-and-asymptotes}
\author{}
\date{}

\begin{document}
\maketitle


For the following exercises, examine the graphs. Identify where the vertical asymptotes are located.

\begin{enumerate}[start=1]
\item \textit{[Figure: The function graphed decreases very rapidly as it approaches $x = 1$ from the left, and on the other side of $x = 1$, it seems to start near infinity and then decrease rapidly.]}
\item \textit{[Figure: The function graphed increases very rapidly as it approaches $x = -3$ from the left, and on the other side of $x = -3$, it seems to start near negative infinity and then increase rapidly to form a sort of U shape that is pointing down, with the other side of the U being at $x = 2$. On the other side of $x = 2$, the graph seems to start near infinity and then decrease rapidly.]}
\item \textit{[Figure: The function graphed decreases very rapidly as it approaches $x = -1$ from the left, and on the other side of $x = -1$, it seems to start near negative infinity and then increase rapidly to form a sort of U shape that is pointing down, with the other side of the U being at $x = 2$. On the other side of $x = 2$, the graph seems to start near infinity and then decrease rapidly.]}
\item \textit{[Figure: The function graphed decreases very rapidly as it approaches $x = 0$ from the left, and on the other side of $x = 0$, it seems to start near infinity and then decrease rapidly to form a sort of U shape that is pointing up, with the other side of the U being at $x = 1$. On the other side of $x = 1$, there is another U shape pointing down, with its other side being at $x = 2$. On the other side of $x = 2$, the graph seems to start near negative infinity and then increase rapidly.]}
\item \textit{[Figure: The function graphed decreases very rapidly as it approaches $x = 0$ from the left, and on the other side of $x = 0$, it seems to start near infinity and then decrease rapidly to form a sort of U shape that is pointing up, with the other side being a normal function that appears as if it will take the entirety of the values of the $x$-axis.]}
\end{enumerate}

For the following functions $f(x)$, determine whether there is an asymptote at $x=a$. Justify your answer without graphing on a calculator.

\begin{enumerate}[resume]
\item $f(x)=\frac{x+1}{x^{2}+5x+4},a=-1$
\item $f(x)=\frac{x}{x-2},a=2$
\item $f(x)={(x+2)}^{3/2},a=-2$
\item $f(x)={(x-1)}^{-1/3},a=1$
\item $f(x)=1+x^{-2/5},a=1$
\end{enumerate}

For the following exercises, evaluate the limit.

\begin{enumerate}[resume]
\item $\lim_{x\to \infty}\frac{1}{3x+6}$
\item $\lim_{x\to \infty}\frac{2x-5}{4x}$
\item $\lim_{x\to \infty}\frac{x^{2}-2x+5}{x+2}$
\item $\lim_{x\to-\infty}\frac{3x^{3}-2x}{x^{2}+2x+8}$
\item $\lim_{x\to-\infty}\frac{x^{4}-4x^{3}+1}{2-2x^{2}-7x^{4}}$
\item $\lim_{x\to \infty}\frac{3x}{\sqrt{x^{2}+1}}$
\item $\lim_{x\to-\infty}\frac{\sqrt{4x^{2}-1}}{x+2}$
\item $\lim_{x\to \infty}\frac{4x}{\sqrt{x^{2}-1}}$
\item $\lim_{x\to-\infty}\frac{4x}{\sqrt{x^{2}-1}}$
\item $\lim_{x\to \infty}\frac{2\sqrt{x}}{x-\sqrt{x}+1}$
\end{enumerate}

For the following exercises, find the horizontal and vertical asymptotes.

\begin{enumerate}[resume]
\item $f(x)=x-\frac{9}{x}$
\item $f(x)=\frac{1}{1-x^{2}}$
\item $f(x)=\frac{x^{3}}{4-x^{2}}$
\item $f(x)=\frac{x^{2}+3}{x^{2}+1}$
\item $f(x)=\sin(x)\sin(2x)$
\item $f(x)=\cos x+\cos(3x)+\cos(5x)$
\item $f(x)=\frac{x\sin(x)}{x^{2}-1}$
\item $f(x)=\frac{x}{\sin(x)}$
\item $f(x)=\frac{1}{x^{3}+x^{2}}$
\item $f(x)=\frac{1}{x-1}-2x$
\item $f(x)=\frac{x^{3}+1}{x^{3}-1}$
\item $f(x)=\frac{\sin x+\cos x}{\sin x-\cos x}$
\item $f(x)=x-\sin x$
\item $f(x)=\frac{1}{x}-\sqrt{x}$
\end{enumerate}

For the following exercises, construct a function $f(x)$ that has the given asymptotes.

\begin{enumerate}[resume]
\item $x=1$ and $y=2$
\item $x=1$ and $y=0$
\item $y=4$, $x=-1$
\item $x=0$
\end{enumerate}

For the following exercises, graph the function on a graphing calculator on the window $x=[-5,5]$ and estimate the horizontal asymptote or limit. Then, calculate the actual horizontal asymptote or limit.

\begin{enumerate}[resume]
\item \techrequired $f(x)=\frac{1}{x+10}$
\item \techrequired $f(x)=\frac{x+1}{x^{2}+7x+6}$
\item \techrequired $\lim_{x\to-\infty}x^{2}+10x+25$
\item \techrequired $\lim_{x\to-\infty}\frac{x+2}{x^{2}+7x+6}$
\item \techrequired $\lim_{x\to \infty}\frac{3x+2}{x+5}$
\end{enumerate}

For the following exercises, draw a graph of the functions without using a calculator. Be sure to notice all important features of the graph: local maxima and minima, inflection points, and asymptotic behavior.

\begin{enumerate}[resume]
\item $y=3x^{2}+2x+4$
\item $y=x^{3}-3x^{2}+4$
\item $y=\frac{2x+1}{x^{2}+6x+5}$
\item $y=\frac{x^{3}+4x^{2}+3x}{3x+9}$
\item $y=\frac{x^{2}+x-2}{x^{2}-3x-4}$
\item $y=\sqrt{x^{2}-5x+4}$
\item $y=2x\sqrt{16-x^{2}}$
\item $y=\frac{\cos x}{x}$, on $x=[-2\pi,2\pi]$
\item $y=\frac{\sqrt{x^{2}+2}}{x+1}$
\item $y=x\tan x,x=[-\pi,\pi]$
\item $y=x\ln(x),x>0$
\item $y=x^{2}\sin(x),x=[-2\pi,2\pi]$
\item For $f(x)=\frac{P(x)}{Q(x)}$ to have an asymptote at $y=2$ then the polynomials $P(x)$ and $Q(x)$ must have what relation?
\item For $f(x)=\frac{P(x)}{Q(x)}$ to have an asymptote at $x=0$, then the polynomials $P(x)$ and $Q(x)$. must have what relation?
\item If $f'(x)$ has asymptotes at $y=3$ and $x=1$, then $f(x)$ has what asymptotes?
\item Both $f(x)=\frac{1}{(x-1)}$ and $g(x)=\frac{1}{{(x-1)}^{2}}$ have asymptotes at $x=1$ and $y=0$. What is the most obvious difference between these two functions?
\item True or false: Every ratio of polynomials has vertical asymptotes.
\end{enumerate}

\end{document}
